\begin{solapartat}
La segona llei de Newton estableix que: $\vec{F} = m\vec{a}$ \quad\punts{0,1}

Segons la llei de gravitació universal, el mòdul de la força s'expressa com:

$F = G\,\dfrac{mM_T}{r^2}$ \quad\punts{0,2}

Per tant, obtenim que: $a = G\,\dfrac{M_T}{r^2}$ \quad\punts{0,1}

D'altra banda, considerant que el satèl·lit descriu un moviment circular uniforme al voltant de la terra, la seva acceleració centrípeta és: $a = \dfrac{v^2}{r}$ o $a = \omega^2 r$ \punts{0,1}\unskip, i la velocitat es pot expressar com $v = \dfrac{2\pi r}{T}$ o $\omega = \dfrac{2\pi}{T}$ \quad\punts{0,1}

Per tant: $T = 2\pi\sqrt{\dfrac{r^3}{GM_T}}$ \quad\punts{0,1}

D'altra banda, $r = (530 + 6370)\times 1000 = 6,9\times 10^{6}\ \mathrm{m}$ \quad\punts{0,1}

Llavors: $T = 2\pi\sqrt{\dfrac{(6,9\times 10^{6})^3}{6,67\times 10^{-11}\cdot 5,98\times 10^{24}}} = 5700\ \mathrm{s}$ \quad\punts{0,1}

I el nombre de voltes és $\dfrac{86400}{5700} = 15,15$, voltes senceres $= 15$ \quad\punts{0,1}
\end{solapartat}

\begin{solapartat}
Intensitat del camp gravitatori a l'òrbita del satèl·lit $g = G\,\dfrac{M_T}{r^2} = 8,38\ \mathrm{m/s^2}$ \quad\punts{0,5}

$\mathit{Pes} = mg = 10,9\ \mathrm{N}$ \quad\punts{0,5}
\end{solapartat}
